% Native editable Beamer slides. Build from the repository root.
\documentclass[aspectratio=169,11pt,handout,t]{beamer}
\makeatletter
\def\input@path{{../}}
\makeatother
\usepackage{satnetslides}
\usepackage{listings}
\definecolor{Icon0}{HTML}{555555}
\definecolor{Icon1}{HTML}{999999}
\definecolor{Icon2}{HTML}{111111}
\definecolor{Icon3}{HTML}{F3F3F3}
\definecolor{Icon4}{HTML}{BDBDBD}
\definecolor{Icon5}{HTML}{666666}
\definecolor{Icon6}{HTML}{F8F8F8}
\definecolor{Icon7}{HTML}{B8B8B8}
\definecolor{Icon8}{HTML}{C5C5C5}
\definecolor{Icon9}{HTML}{DEDEDE}
\definecolor{Icon10}{HTML}{A8A8A8}
\definecolor{Icon11}{HTML}{FAFAFA}
\definecolor{Icon12}{HTML}{8E8E8E}
\definecolor{Icon13}{HTML}{E6E6E6}
\definecolor{Icon14}{HTML}{CCCCCC}
\definecolor{Icon15}{HTML}{E2E2E2}
\definecolor{Icon16}{HTML}{D5D5D5}
\definecolor{Icon17}{HTML}{F7F7F7}
\definecolor{Icon18}{HTML}{DADADA}
\definecolor{Icon19}{HTML}{777777}
\definecolor{Icon20}{HTML}{DDDDDD}
\definecolor{Icon21}{HTML}{C8C8C8}
\definecolor{Icon22}{HTML}{F0F0F0}
\definecolor{Icon23}{HTML}{CFCFCF}
\definecolor{Icon24}{HTML}{C6C6C6}
\definecolor{Icon25}{HTML}{E0E0E0}
\definecolor{Icon26}{HTML}{E5E5E5}
\definecolor{Icon27}{HTML}{4B4B4B}
\definecolor{Icon28}{HTML}{FFFFFF}
\definecolor{Icon29}{HTML}{BEBEBE}
\definecolor{Icon30}{HTML}{F4F4F4}
\definecolor{Icon31}{HTML}{ECECEC}
% Native TikZ paths for the shared monochrome component icons.
% Both solar arrays use the same axonometric projection.
\tikzset{pics/satellite/.style={code={%
\path[fill=none,draw=Icon0,line width=1.14pt,line join=round,line cap=round] (-0.276,-0.00092) -- (0.276,0.12052);
\path[fill=Icon1,draw=Icon2,line width=0.304pt,line join=round,line cap=round] (-0.5244,-0.15346) -- (-0.184,-0.07857) -- (-0.184,-0.09697) -- (-0.5244,-0.17186) -- cycle;
\path[fill=Icon3,draw=Icon2,line width=0.494pt,line join=round,line cap=round] (-0.6532,0.01398) -- (-0.3128,0.08887) -- (-0.184,-0.07857) -- (-0.5244,-0.15346) -- cycle;
\path[fill=Icon4,draw=Icon5,line width=0.19pt,line join=round,line cap=round] (-0.62836,0.00546) -- (-0.31924,0.07347) -- (-0.20884,-0.07005) -- (-0.51796,-0.13806) -- cycle;
\path[fill=none,draw=Icon6,line width=0.247pt,line join=round,line cap=round] (-0.57684,0.0168) -- (-0.46644,-0.12672);
\path[fill=none,draw=Icon6,line width=0.247pt,line join=round,line cap=round] (-0.52532,0.02813) -- (-0.41492,-0.11539);
\path[fill=none,draw=Icon6,line width=0.247pt,line join=round,line cap=round] (-0.4738,0.03947) -- (-0.3634,-0.10405);
\path[fill=none,draw=Icon6,line width=0.247pt,line join=round,line cap=round] (-0.42228,0.0508) -- (-0.31188,-0.09272);
\path[fill=none,draw=Icon6,line width=0.247pt,line join=round,line cap=round] (-0.37076,0.06214) -- (-0.26036,-0.08138);
\path[fill=none,draw=Icon6,line width=0.247pt,line join=round,line cap=round] (-0.59156,-0.04238) -- (-0.28244,0.02563);
\path[fill=none,draw=Icon6,line width=0.247pt,line join=round,line cap=round] (-0.55476,-0.09022) -- (-0.24564,-0.02221);
\path[fill=none,draw=Icon5,line width=0.304pt,line join=round,line cap=round] (-0.483,0.05143) -- (-0.3542,-0.11601);
\path[fill=Icon1,draw=Icon2,line width=0.304pt,line join=round,line cap=round] (0.3128,0.03073) -- (0.6532,0.10562) -- (0.6532,0.08722) -- (0.3128,0.01233) -- cycle;
\path[fill=Icon3,draw=Icon2,line width=0.494pt,line join=round,line cap=round] (0.184,0.19817) -- (0.5244,0.27306) -- (0.6532,0.10562) -- (0.3128,0.03073) -- cycle;
\path[fill=Icon4,draw=Icon5,line width=0.19pt,line join=round,line cap=round] (0.20884,0.18965) -- (0.51796,0.25766) -- (0.62836,0.11414) -- (0.31924,0.04613) -- cycle;
\path[fill=none,draw=Icon6,line width=0.247pt,line join=round,line cap=round] (0.26036,0.20098) -- (0.37076,0.05746);
\path[fill=none,draw=Icon6,line width=0.247pt,line join=round,line cap=round] (0.31188,0.21232) -- (0.42228,0.0688);
\path[fill=none,draw=Icon6,line width=0.247pt,line join=round,line cap=round] (0.3634,0.22365) -- (0.4738,0.08013);
\path[fill=none,draw=Icon6,line width=0.247pt,line join=round,line cap=round] (0.41492,0.23499) -- (0.52532,0.09147);
\path[fill=none,draw=Icon6,line width=0.247pt,line join=round,line cap=round] (0.46644,0.24632) -- (0.57684,0.1028);
\path[fill=none,draw=Icon6,line width=0.247pt,line join=round,line cap=round] (0.24564,0.14181) -- (0.55476,0.20982);
\path[fill=none,draw=Icon6,line width=0.247pt,line join=round,line cap=round] (0.28244,0.09397) -- (0.59156,0.16198);
\path[fill=none,draw=Icon5,line width=0.304pt,line join=round,line cap=round] (0.3542,0.23561) -- (0.483,0.06817);
\path[fill=Icon7,draw=Icon2,line width=0.418pt,line join=round,line cap=round] (-0.1978,0.18961) -- (-0.115,0.08197) -- (-0.115,-0.07443) -- (-0.1978,0.03321) -- cycle;
\path[fill=Icon8,draw=Icon2,line width=0.418pt,line join=round,line cap=round] (-0.115,0.08197) -- (-0.0736,0.0701) -- (-0.0736,-0.0863) -- (-0.115,-0.07443) -- cycle;
\path[fill=Icon9,draw=Icon2,line width=0.418pt,line join=round,line cap=round] (-0.0736,0.0701) -- (0.184,0.12678) -- (0.184,-0.02962) -- (-0.0736,-0.0863) -- cycle;
\path[fill=Icon7,draw=Icon2,line width=0.418pt,line join=round,line cap=round] (0.184,0.12678) -- (0.1978,0.15079) -- (0.1978,-0.00561) -- (0.184,-0.02962) -- cycle;
\path[fill=Icon10,draw=Icon2,line width=0.418pt,line join=round,line cap=round] (0.1978,0.15079) -- (0.115,0.25843) -- (0.115,0.10203) -- (0.1978,-0.00561) -- cycle;
\path[fill=Icon11,draw=Icon2,line width=0.475pt,line join=round,line cap=round] (-0.184,0.21362) -- (0.0736,0.2703) -- (0.115,0.25843) -- (0.1978,0.15079) -- (0.184,0.12678) -- (-0.0736,0.0701) -- (-0.115,0.08197) -- (-0.1978,0.18961) -- cycle;
\path[fill=Icon12,draw=Icon2,line width=0.266pt,line join=round,line cap=round] (-0.0276,0.04342) -- (0.1472,0.08188) -- (0.1472,-0.00092) -- (-0.0276,-0.03938) -- cycle;
\path[fill=none,draw=Icon13,line width=0.209pt,line join=round,line cap=round] (0.0161,0.04384) -- (0.0161,-0.02056);
\path[fill=none,draw=Icon13,line width=0.209pt,line join=round,line cap=round] (0.0598,0.05345) -- (0.0598,-0.01095);
\path[fill=none,draw=Icon13,line width=0.209pt,line join=round,line cap=round] (0.1035,0.06307) -- (0.1035,-0.00133);
\path[fill=Icon14,draw=Icon2,line width=0.304pt,line join=round,line cap=round] (0.0782,0.20599) -- (0.08144,0.20056) -- (0.08297,0.195) -- (0.08272,0.1895) -- (0.08071,0.18429) -- (0.07701,0.17957) -- (0.07177,0.1755) -- (0.06518,0.17226) -- (0.0575,0.16997) -- (0.04903,0.16871) -- (0.04008,0.16853) -- (0.03101,0.16945) -- (0.02216,0.17141) -- (0.01388,0.17436) -- (0.00647,0.17817) -- (0.00023,0.1827) -- (-0.0046,0.18777) -- (-0.00784,0.1932) -- (-0.00937,0.19876) -- (-0.00912,0.20426) -- (-0.00711,0.20947) -- (-0.00341,0.21419) -- (0.00183,0.21826) -- (0.00842,0.2215) -- (0.0161,0.22379) -- (0.02457,0.22505) -- (0.03352,0.22523) -- (0.04259,0.22431) -- (0.05144,0.22235) -- (0.05972,0.2194) -- (0.06713,0.21559) -- (0.07337,0.21106) -- cycle;
\path[fill=none,draw=Icon2,line width=0.437pt,line join=round,line cap=round] (-0.1012,0.1759) -- (-0.1012,0.2955);
\path[fill=none,draw=Icon2,line width=0.38pt,line join=round,line cap=round] (-0.1288,0.28943) -- (-0.0736,0.30158);
}}}
\tikzset{pics/user/.style={code={%
\path[fill=Icon15,draw=Icon2,line width=0.494pt,line join=round,line cap=round] (-0.2408,0.172) ellipse [x radius=0.0688cm,y radius=0.0688cm];
\path[fill=Icon16,draw=Icon2,line width=0.532pt,line join=round,line cap=round] (-0.3612,-0.0946) -- (-0.3526,0.0086) -- (-0.3096,0.0602) -- (-0.1806,0.0602) -- (-0.1204,0.0086) -- (-0.086,-0.0946) -- cycle;
\path[fill=Icon17,draw=Icon2,line width=0.532pt,line join=round,line cap=round] (-0.1118,0.0946) -- (0.215,0.0946) -- (0.1806,-0.1376) -- (-0.1376,-0.1376) -- cycle;
\path[fill=Icon18,draw=Icon19,line width=0.304pt,line join=round,line cap=round] (-0.0774,0.0602) -- (0.172,0.0602) -- (0.1462,-0.0946) -- (-0.1032,-0.0946) -- cycle;
\path[fill=Icon20,draw=Icon2,line width=0.532pt,line join=round,line cap=round] (-0.1376,-0.1376) -- (0.1806,-0.1376) -- (0.2752,-0.1978) -- (-0.2408,-0.1978) -- cycle;
\path[fill=none,draw=Icon19,line width=0.323pt,line join=round,line cap=round] (-0.1032,-0.1634) -- (0.1032,-0.1634);
\path[fill=none,draw=Icon2,line width=0.57pt,line join=round,line cap=round] (-0.3698,-0.2322) -- (0.3096,-0.2322);
}}}
\tikzset{pics/user-router/.style={code={%
\path[fill=Icon15,draw=Icon2,line width=0.494pt,line join=round,line cap=round] (-0.2408,0.172) ellipse [x radius=0.0688cm,y radius=0.0688cm];
\path[fill=Icon16,draw=Icon2,line width=0.532pt,line join=round,line cap=round] (-0.3612,-0.0946) -- (-0.3526,0.0086) -- (-0.3096,0.0602) -- (-0.1806,0.0602) -- (-0.1204,0.0086) -- (-0.086,-0.0946) -- cycle;
\path[fill=Icon17,draw=Icon2,line width=0.532pt,line join=round,line cap=round] (-0.1118,0.0946) -- (0.215,0.0946) -- (0.1806,-0.1376) -- (-0.1376,-0.1376) -- cycle;
\path[fill=Icon18,draw=Icon19,line width=0.304pt,line join=round,line cap=round] (-0.0774,0.0602) -- (0.172,0.0602) -- (0.1462,-0.0946) -- (-0.1032,-0.0946) -- cycle;
\path[fill=Icon20,draw=Icon2,line width=0.532pt,line join=round,line cap=round] (-0.1376,-0.1376) -- (0.1806,-0.1376) -- (0.2752,-0.1978) -- (-0.2408,-0.1978) -- cycle;
\path[fill=none,draw=Icon19,line width=0.323pt,line join=round,line cap=round] (-0.1032,-0.1634) -- (0.1032,-0.1634);
\path[fill=none,draw=Icon2,line width=0.57pt,line join=round,line cap=round] (-0.3698,-0.2322) -- (0.3096,-0.2322);
\path[fill=Icon21,draw=Icon2,line width=0.494pt,line join=round,line cap=round] (0.2322,-0.0258) -- (0.3956,-0.0258) -- (0.3956,-0.1376) -- (0.2322,-0.1376) -- cycle;
\path[fill=none,draw=Icon2,line width=0.57pt,line join=round,line cap=round] (0.344,-0.0258) -- (0.344,0.1032);
\path[fill=Icon2,draw=Icon2,line width=0.152pt,line join=round,line cap=round] (0.27348,-0.08428) ellipse [x radius=0.00688cm,y radius=0.00688cm];
\path[fill=Icon2,draw=Icon2,line width=0.152pt,line join=round,line cap=round] (0.31648,-0.08428) ellipse [x radius=0.00688cm,y radius=0.00688cm];
\path[fill=Icon2,draw=Icon2,line width=0.152pt,line join=round,line cap=round] (0.35948,-0.08428) ellipse [x radius=0.00688cm,y radius=0.00688cm];
}}}
\tikzset{pics/terminal/.style={code={%
\path[fill=Icon4,draw=Icon2,line width=0.532pt,line join=round,line cap=round] (-0.215,0.1892) -- (0.258,0.1204) -- (0.215,-0.043) -- (-0.258,0.0258) -- cycle;
\path[fill=Icon22,draw=Icon2,line width=0.532pt,line join=round,line cap=round] (-0.215,0.215) -- (0.258,0.1462) -- (0.2236,-0.0086) -- (-0.2494,0.0602) -- cycle;
\path[fill=none,draw=Icon1,line width=0.285pt,line join=round,line cap=round] (-0.1806,0.1806) -- (0.215,0.1204);
\path[fill=none,draw=Icon0,line width=1.14pt,line join=round,line cap=round] (-0.0086,0.0172) -- (-0.0086,-0.1376);
\path[fill=Icon23,draw=Icon2,line width=0.532pt,line join=round,line cap=round] (-0.0344,-0.1204) -- (0.0172,-0.1204) -- (0.1204,-0.2408) -- (0.0516,-0.2408) -- (-0.0086,-0.1634) -- (-0.1376,-0.2408) -- (-0.215,-0.2408) -- cycle;
\path[fill=none,draw=Icon2,line width=0.76pt,line join=round,line cap=round] (-0.0086,-0.1376) -- (0.1548,-0.1806);
}}}
\tikzset{pics/gateway/.style={code={%
\path[fill=Icon24,draw=Icon2,line width=0.532pt,line join=round,line cap=round] (-0.1462,-0.1978) -- (0.1376,-0.1978) -- (0.1978,-0.258) -- (-0.215,-0.258) -- cycle;
\path[fill=Icon16,draw=Icon2,line width=0.532pt,line join=round,line cap=round] (-0.0258,-0.0258) -- (0.0516,-0.0344) -- (0.0946,-0.1978) -- (-0.086,-0.1978) -- cycle;
\path[fill=Icon25,draw=Icon2,line width=0.532pt,line join=round,line cap=round] (-0.25717,0.03102) -- (-0.23001,0.00839) -- (-0.20385,-0.0121) -- (-0.1787,-0.03043) -- (-0.15454,-0.04662) -- (-0.13139,-0.06066) -- (-0.10924,-0.07255) -- (-0.08809,-0.08229) -- (-0.06794,-0.08989) -- (-0.04879,-0.09534) -- (-0.03065,-0.09865) -- (-0.0135,-0.0998) -- (0.00264,-0.09881) -- (0.01778,-0.09567) -- (0.03192,-0.09038) -- (0.04506,-0.08295) -- (0.0572,-0.07337) -- (0.06834,-0.06164) -- (0.07847,-0.04776) -- (0.08761,-0.03174) -- (0.09574,-0.01357) -- (0.10287,0.00675) -- (0.109,0.02922) -- (0.11413,0.05383) -- (0.11825,0.08059) -- (0.12138,0.1095) -- (0.12351,0.14055) -- (0.12463,0.17376) -- (0.12475,0.20911) -- cycle;
\path[fill=Icon11,draw=Icon2,line width=0.532pt,line join=round,line cap=round] (0.12475,0.20911) -- (0.12581,0.2007) -- (0.12214,0.1903) -- (0.11384,0.17817) -- (0.1011,0.16461) -- (0.08424,0.14996) -- (0.06368,0.13457) -- (0.03991,0.11882) -- (0.01354,0.1031) -- (-0.0148,0.0878) -- (-0.0444,0.0733) -- (-0.07454,0.05994) -- (-0.10448,0.04807) -- (-0.13347,0.03797) -- (-0.16081,0.02989) -- (-0.18582,0.02403) -- (-0.20788,0.02053) -- (-0.22645,0.01949) -- (-0.24108,0.02092) -- (-0.25141,0.0248) -- (-0.25717,0.03102) -- (-0.25823,0.03943) -- (-0.25456,0.04983) -- (-0.24626,0.06195) -- (-0.23352,0.07551) -- (-0.21666,0.09017) -- (-0.19609,0.10556) -- (-0.17233,0.1213) -- (-0.14596,0.13702) -- (-0.11762,0.15232) -- (-0.08802,0.16683) -- (-0.05788,0.18018) -- (-0.02794,0.19206) -- (0.00105,0.20216) -- (0.02839,0.21024) -- (0.0534,0.21609) -- (0.07546,0.21959) -- (0.09404,0.22063) -- (0.10866,0.2192) -- (0.11899,0.21533) -- cycle;
\path[fill=none,draw=Icon2,line width=0.437pt,line join=round,line cap=round] (-0.24158,0.03829) -- (-0.11346,0.22139);
\path[fill=none,draw=Icon2,line width=0.437pt,line join=round,line cap=round] (0.10916,0.20184) -- (-0.11346,0.22139);
\path[fill=Icon5,draw=Icon2,line width=0.494pt,line join=round,line cap=round] (-0.11346,0.22139) ellipse [x radius=0.0215cm,y radius=0.0215cm];
\path[fill=Icon26,draw=Icon2,line width=0.494pt,line join=round,line cap=round] (0.2322,-0.086) -- (0.3784,-0.086) -- (0.3784,-0.2494) -- (0.2322,-0.2494) -- cycle;
\path[fill=none,draw=Icon19,line width=0.323pt,line join=round,line cap=round] (0.258,-0.129) -- (0.344,-0.129);
\path[fill=none,draw=Icon19,line width=0.323pt,line join=round,line cap=round] (0.258,-0.1634) -- (0.344,-0.1634);
\path[fill=none,draw=Icon19,line width=0.323pt,line join=round,line cap=round] (0.258,-0.1978) -- (0.344,-0.1978);
}}}
\tikzset{pics/internet/.style={code={%
\path[fill=Icon17,draw=Icon2,line width=0.494pt,line join=round,line cap=round] (0,0.0086) ellipse [x radius=0.2408cm,y radius=0.2408cm];
\path[fill=none,draw=Icon19,line width=0.38pt,line join=round,line cap=round] (0,0.0086) ellipse [x radius=0.129cm,y radius=0.2408cm];
\path[fill=none,draw=Icon19,line width=0.38pt,line join=round,line cap=round] (0,0.0086) ellipse [x radius=0.2408cm,y radius=0.0946cm];
\path[fill=none,draw=Icon1,line width=0.342pt,line join=round,line cap=round] (-0.2408,0.0086) -- (0.2408,0.0086);
\path[fill=none,draw=Icon1,line width=0.342pt,line join=round,line cap=round] (0,0.2494) -- (0,-0.2322);
\path[fill=none,draw=Icon2,line width=0.475pt,line join=round,line cap=round] (-0.1806,0.1204) -- (0.1032,0.172);
\path[fill=none,draw=Icon2,line width=0.475pt,line join=round,line cap=round] (0.1032,0.172) -- (0.172,-0.0946);
\path[fill=none,draw=Icon2,line width=0.475pt,line join=round,line cap=round] (0.172,-0.0946) -- (-0.1204,-0.1462);
\path[fill=none,draw=Icon2,line width=0.475pt,line join=round,line cap=round] (-0.1204,-0.1462) -- (-0.1806,0.1204);
\path[fill=Icon27,draw=Icon28,line width=0.304pt,line join=round,line cap=round] (-0.1806,0.1204) ellipse [x radius=0.0258cm,y radius=0.0258cm];
\path[fill=Icon27,draw=Icon28,line width=0.304pt,line join=round,line cap=round] (0.1032,0.172) ellipse [x radius=0.0258cm,y radius=0.0258cm];
\path[fill=Icon27,draw=Icon28,line width=0.304pt,line join=round,line cap=round] (0.172,-0.0946) ellipse [x radius=0.0258cm,y radius=0.0258cm];
\path[fill=Icon27,draw=Icon28,line width=0.304pt,line join=round,line cap=round] (-0.1204,-0.1462) ellipse [x radius=0.0258cm,y radius=0.0258cm];
}}}
\tikzset{pics/server/.style={code={%
\path[fill=Icon29,draw=Icon2,line width=0.532pt,line join=round,line cap=round] (0.1462,0.2064) -- (0.2494,0.1376) -- (0.2494,-0.258) -- (0.1462,-0.1978) -- cycle;
\path[fill=Icon30,draw=Icon2,line width=0.532pt,line join=round,line cap=round] (-0.2064,0.2064) -- (0.1462,0.2064) -- (0.2494,0.1376) -- (-0.1032,0.1376) -- cycle;
\path[fill=Icon31,draw=Icon2,line width=0.494pt,line join=round,line cap=round] (-0.2064,0.2064) -- (0.1462,0.2064) -- (0.1462,-0.215) -- (-0.2064,-0.215) -- cycle;
\path[fill=Icon28,draw=Icon0,line width=0.38pt,line join=round,line cap=round] (-0.172,0.1548) -- (0.1118,0.1548) -- (0.1118,0.0688) -- (-0.172,0.0688) -- cycle;
\path[fill=Icon5,draw=none,line width=0pt,line join=round,line cap=round] (-0.13072,0.1118) ellipse [x radius=0.01548cm,y radius=0.01548cm];
\path[fill=none,draw=Icon19,line width=0.304pt,line join=round,line cap=round] (-0.0602,0.129) -- (-0.0602,0.0946);
\path[fill=none,draw=Icon19,line width=0.304pt,line join=round,line cap=round] (-0.0258,0.129) -- (-0.0258,0.0946);
\path[fill=none,draw=Icon19,line width=0.304pt,line join=round,line cap=round] (0.0086,0.129) -- (0.0086,0.0946);
\path[fill=none,draw=Icon19,line width=0.304pt,line join=round,line cap=round] (0.043,0.129) -- (0.043,0.0946);
\path[fill=Icon28,draw=Icon0,line width=0.38pt,line join=round,line cap=round] (-0.172,0.043) -- (0.1118,0.043) -- (0.1118,-0.043) -- (-0.172,-0.043) -- cycle;
\path[fill=Icon5,draw=none,line width=0pt,line join=round,line cap=round] (-0.13072,0) ellipse [x radius=0.01548cm,y radius=0.01548cm];
\path[fill=none,draw=Icon19,line width=0.304pt,line join=round,line cap=round] (-0.0602,0.0172) -- (-0.0602,-0.0172);
\path[fill=none,draw=Icon19,line width=0.304pt,line join=round,line cap=round] (-0.0258,0.0172) -- (-0.0258,-0.0172);
\path[fill=none,draw=Icon19,line width=0.304pt,line join=round,line cap=round] (0.0086,0.0172) -- (0.0086,-0.0172);
\path[fill=none,draw=Icon19,line width=0.304pt,line join=round,line cap=round] (0.043,0.0172) -- (0.043,-0.0172);
\path[fill=Icon28,draw=Icon0,line width=0.38pt,line join=round,line cap=round] (-0.172,-0.0688) -- (0.1118,-0.0688) -- (0.1118,-0.1548) -- (-0.172,-0.1548) -- cycle;
\path[fill=Icon5,draw=none,line width=0pt,line join=round,line cap=round] (-0.13072,-0.1118) ellipse [x radius=0.01548cm,y radius=0.01548cm];
\path[fill=none,draw=Icon19,line width=0.304pt,line join=round,line cap=round] (-0.0602,-0.0946) -- (-0.0602,-0.129);
\path[fill=none,draw=Icon19,line width=0.304pt,line join=round,line cap=round] (-0.0258,-0.0946) -- (-0.0258,-0.129);
\path[fill=none,draw=Icon19,line width=0.304pt,line join=round,line cap=round] (0.0086,-0.0946) -- (0.0086,-0.129);
\path[fill=none,draw=Icon19,line width=0.304pt,line join=round,line cap=round] (0.043,-0.0946) -- (0.043,-0.129);
\path[fill=none,draw=Icon2,line width=0.874pt,line join=round,line cap=round] (-0.1806,-0.2322) -- (-0.1032,-0.2322);
\path[fill=none,draw=Icon2,line width=0.874pt,line join=round,line cap=round] (0.0516,-0.2322) -- (0.129,-0.2322);
}}}

\lstset{language=Python,basicstyle=\ttfamily\fontsize{7.5}{9.5}\selectfont,
  keywordstyle=\bfseries\color{Ink},commentstyle=\color{Muted},
  stringstyle=\color{Muted},showstringspaces=false,columns=fullflexible,
  keepspaces=true,frame=single,rulecolor=\color{Line},xleftmargin=3pt,xrightmargin=3pt}
\tikzset{box/.style={draw=Line,fill=Soft,rounded corners=1pt,font=\scriptsize,
  align=center,minimum height=0.65cm,inner sep=4pt},
  dot/.style={circle,draw=Ink,fill=Paper,minimum size=0.50cm,inner sep=0pt,font=\scriptsize\bfseries}}
\pgfplotsset{courseaxis/.style={width=6.6cm,height=4.0cm,
  label style={font=\scriptsize},tick label style={font=\scriptsize},
  axis line style={Muted},tick style={Muted},grid=major,grid style={Line,line width=0.25pt},
  legend style={font=\scriptsize,draw=none,fill=Paper},unbounded coords=jump}}
\setlecture{05}{Designing and Evaluating Satellite Networks}
\title{Designing and Evaluating Satellite Networks}
\author{Yi Ching (David) Chou}
\date{}
\begin{document}
\begin{frame}[plain]
  \begin{tikzpicture}[remember picture,overlay]
    \fill[Paper] (current page.south west) rectangle (current page.north east);
    \fill[Line] (current page.south west) rectangle ([yshift=0.18cm]current page.south east);
    \fill[Ink] ([xshift=0.95cm,yshift=-0.95cm]current page.north west) rectangle ++(0.12cm,-5.35cm);
    \node[anchor=north west,text=Muted] at ([xshift=1.35cm,yshift=-0.9cm]current page.north west)
      {\small\bfseries SATNET EDU / Lecture 05};
    \node[anchor=north west,align=left,text width=12.2cm]
      at ([xshift=1.35cm,yshift=-1.65cm]current page.north west)
      {{\fontsize{29}{33}\selectfont\bfseries Designing and Evaluating\\[0.12cm]Satellite Networks}};
    \node[anchor=north west,align=left,text width=11.6cm,text=Muted]
      at ([xshift=1.38cm,yshift=-3.85cm]current page.north west)
      {\small How to turn a network idea into a controlled, reproducible experiment.};
    \node[anchor=north west,fill=Ink,rounded corners=2pt,inner xsep=0.18cm,inner ysep=0.09cm,text=Paper]
      (tag) at ([xshift=1.38cm,yshift=-4.78cm]current page.north west)
      {\small\bfseries LECTURE 05};
    \node[anchor=west,align=left,text width=9.5cm] at ([xshift=0.28cm]tag.east)
      {\small From a design question to an explained result};
    \draw[Ink,line width=1.1pt] ([xshift=1.38cm,yshift=-5.55cm]current page.north west) -- ++(10.9cm,0);
    \node[anchor=north,align=center,text=Muted] at ([yshift=-5.85cm]current page.north)
      {\small State the question. Control the comparison. Inspect the outcome.};
    \node[anchor=north,align=center] at ([yshift=-6.65cm]current page.north)
      {\small Instructor: Yi Ching (David) Chou};
  \end{tikzpicture}
\end{frame}

\begin{frame}[fragile]{Start with a Question That the Model Can Answer}
  \context{Connect a design choice to a measurable outcome under stated conditions.}
  \begin{columns}[T,onlytextwidth]
    \begin{column}{0.47\textwidth}\raggedright
      \begin{teachingpoints}
        \point{Turn a goal into a comparison.}{\item Ask how a specific topology change affects connectivity between two chosen ground sites.}
        \point{Predict a mechanism first.}{\item Allowing more candidate ISLs may create alternate paths. Range and Earth blockage can still remove those links.}
        \point{Match the claim to the model.}{\item SatNet Edu can compare geometric reachability and propagation. It cannot predict throughput without a traffic model.}
      \end{teachingpoints}
    \end{column}
    \begin{column}{0.49\textwidth}\centering
      \figtitle{Make the question testable}
\begin{tikzpicture}[x=0.90cm,y=0.90cm]
\path[use as bounding box] (-0.25,-0.2) rectangle (7.15,4.5);

\node[box,text width=5.5cm] (q) at (3.35,3.9) {Question\\Can an alternate path survive a failure?};
\node[box,text width=5.5cm] (h) at (3.35,2.55) {Prediction\\Avoiding the failed node may preserve access};
\node[box,text width=5.5cm] (m) at (3.35,1.2) {Measurement\\Count reachable samples during the failure};
\draw[flow] (q)--(h);\draw[flow] (h)--(m);
\node[smalllab,text=Muted] at (3.35,0.15) {Fix endpoints, time window, geometry, and routing metric};

\end{tikzpicture}
    \end{column}
  \end{columns}
  \assumption{The teaching model omits antenna scheduling, capacity, queues, packet loss, and routing convergence.}
  \note{A hypothesis may be rejected. More candidate links need not add a usable alternate path, and a second path sharing the failed node offers no protection. Refer to docs/MODEL.md for the implemented model.}
\end{frame}

\begin{frame}[fragile]{Change One Factor and Keep the Rest Comparable}
  \context{A controlled comparison helps identify why a measured result changed.}
  \begin{columns}[T,onlytextwidth]
    \begin{column}{0.47\textwidth}\raggedright
      \begin{teachingpoints}
        \point{Use a shared baseline.}{\item Use the same epoch, constellation, sites, elevation mask, duration, and sample times in both runs.}
        \point{Name the factor you change.}{\item Here the failure schedule is the treatment. Add one satellite outage while leaving geometry and routing settings fixed.}
        \point{Separate follow-up questions.}{\item A later experiment can change topology or altitude. Changing both at once makes a single-cause explanation harder.}
      \end{teachingpoints}
    \end{column}
    \begin{column}{0.49\textwidth}\centering
      \figtitle{The paired failure experiment}
{\scriptsize\setlength{\tabcolsep}{4pt}\renewcommand{\arraystretch}{1.35}
\begin{tabular}{@{}lll@{}}\toprule
\textbf{Setting} & \textbf{Baseline} & \textbf{Failure run} \\ \midrule
Satellites & $6\times12$ & $6\times12$ \\
Altitude / tilt & 1000 km / $53^\circ$ & 1000 km / $53^\circ$ \\
Access mask & $10^\circ$ & $10^\circ$ \\
ISL limit & 4000 km & 4000 km \\
Window / step & 20 min / 10 s & 20 min / 10 s \\
Outage & None & 600--900 s \\
\bottomrule\end{tabular}\par}\par\vspace{0.12cm}{\footnotesize Vancouver $\longleftrightarrow$ Tokyo\quad Same epoch\par}
    \end{column}
  \end{columns}
  \assumption{Companion experiment uses orbital-neighbor links and epoch 2026-01-01 UTC. The node selection rule is documented.}
  \note{The companion selects an interior satellite on the baseline delay route at 600 seconds and disables that same satellite in the failure run for [600,900). This is a targeted stress case, not a random failure estimate.}
\end{frame}

\begin{frame}[fragile]{Coverage, Reachability, and Delay Need Different Counts}
  \context{Always state the population used in a metric, especially when some samples have no route.}
  \begin{columns}[T,onlytextwidth]
    \begin{column}{0.47\textwidth}\raggedright
      \begin{teachingpoints}
        \point{Coverage is local to one site.}{\item Count samples with at least one usable ground link. This does not require a route to the destination.}
        \point{Reachability is a pairwise property.}{\item Count samples with a complete source-to-destination path, using all scheduled observation times as the denominator.}
        \point{Delay is conditional on a route.}{\item Average valid delays only and report the reachability fraction alongside them. Missing paths are not zero-delay paths.}
      \end{teachingpoints}
    \end{column}
    \begin{column}{0.49\textwidth}\centering
      \figtitle{The denominator changes the meaning}
{\scriptsize\setlength{\tabcolsep}{4pt}\renewcommand{\arraystretch}{1.35}
\begin{tabular}{@{}llll@{}}\toprule
\textbf{Time} & \textbf{U access} & \textbf{U--V route} & \textbf{Delay} \\ \midrule
0 s & Yes & Yes & 12 ms \\
10 s & Yes & No & --- \\
20 s & Yes & Yes & 18 ms \\
30 s & No & No & --- \\
\bottomrule\end{tabular}\par}\par\vspace{0.12cm}{\footnotesize Coverage at U: $3/4=75\%$\par}\par\vspace{0.12cm}{\footnotesize Pair reachability: $2/4=50\%$\par}\par\vspace{0.12cm}{\footnotesize Conditional delay: $(12+18)/2=15$ ms\par}
    \end{column}
  \end{columns}
  \assumption{Constructed four-sample example. Delay here means one-way propagation. These are sample fractions, not exact time fractions.}
  \note{A denominator of four with zeros substituted for missing delays gives 7.5 ms, a misleading score. A mean of 15 ms is valid only when paired with the statement that half the samples had no route.}
\end{frame}

\begin{frame}[fragile]{Unequal Sample Spacing Requires Time Weights}
  \context{Extra samples near an event can change a sample average without changing the time history.}
  \begin{columns}[T,onlytextwidth]
    \begin{column}{0.47\textwidth}\raggedright
      \begin{teachingpoints}
        \point{Samples can represent unequal time.}{\item A recorder can add samples at failure times. Counting every observation equally can overweight short periods.}
        \point{Choose and label an estimator.}{\item A left-hold estimate assumes the observed state lasts until the next sample. Weight each state by that interval length.}
        \point{Check sensitivity to finer sampling.}{\item This assumption can miss an unobserved change. Report the interval, window, and approximation with the metric.}
      \end{teachingpoints}
    \end{column}
    \begin{column}{0.49\textwidth}\centering
      \figtitle{Two reachable samples, unequal durations}
\begin{tikzpicture}[x=0.90cm,y=0.90cm]
\path[use as bounding box] (-0.25,-0.2) rectangle (7.15,4.5);

\draw[flow] (0.65,2.9)--(6.7,2.9);
\draw[Ink,line width=3pt] (0.8,3.1)--(1.32,3.1);
\draw[Ink,line width=3pt] (6.0,3.1)--(6.52,3.1);
\foreach \x/\t/\state in {0.8/0/1,1.32/1/0,6.0/10/1,6.52/11/0}{
 \fill[Ink] (\x,2.9) circle (0.045);
 \node[smalllab,below] at (\x,2.7) {\t};
 \node[smalllab,above] at (\x,3.35) {$\state$};
}
\node[lab] at (3.5,4.05) {$a$: 1 for reachable, 0 otherwise};
\node[lab] at (3.5,1.8) {Sample fraction: $2/4=50\%$};
\node[lab] at (3.5,0.9) {Held-state time fraction: $(1+1)/11=18.2\%$};
\node[smalllab,text=Muted] at (3.5,0.15) {Time in seconds. Final sample has zero following duration.};

\end{tikzpicture}
    \end{column}
  \end{columns}
  \assumption{Constructed states at 0, 1, 10, and 11 s. The time fraction assumes the pictured state is held between samples.}
  \note{The left-hold estimate is sum a\_k times (t\_{k+1}-t\_k), divided by t\_last-t\_first. It is an estimate for a sampled geometric process. An event-aligned schedule can make the held-state assumption exact for that explicit event alone.}
\end{frame}

\begin{frame}[fragile]{A Failure Test Should Target a Defined Component}
  \context{A failure experiment is meaningful only when the disabled component and its timing are specified.}
  \begin{columns}[T,onlytextwidth]
    \begin{column}{0.47\textwidth}\raggedright
      \begin{teachingpoints}
        \point{Choose a repeatable selection rule.}{\item Select an interior satellite on the baseline route at 600 seconds. Save its ID so the stress case can be reproduced.}
        \point{Remove every incident link.}{\item A disabled satellite cannot forward traffic. Removing just one displayed route edge would test a different fault.}
        \point{Look for a surviving alternate path.}{\item The replacement can be longer, or no path may remain. Neither outcome should be assumed before querying the graph.}
      \end{teachingpoints}
    \end{column}
    \begin{column}{0.49\textwidth}\centering
      \figtitle{One failed node affects multiple links}
\begin{tikzpicture}[x=0.90cm,y=0.90cm]
\path[use as bounding box] (-0.25,-0.2) rectangle (7.15,4.5);

\node[dot] (U) at (0.5,2.0) {U};\node[dot] (A) at (2.2,3.4) {A};
\node[dot,fill=Soft] (X) at (3.5,2.0) {X};\node[dot] (B) at (4.8,3.4) {B};
\node[dot] (C) at (3.5,0.65) {C};\node[dot] (V) at (6.5,2.0) {V};
\draw[net] (U)--(A) (B)--(V);
\draw[net,dashed,Muted] (A)--(X)--(B) (X)--(C);
\draw[net,line width=1.1pt] (A)--(C)--(B);
\draw[Ink,line width=1pt] (3.2,1.7)--(3.8,2.3) (3.2,2.3)--(3.8,1.7);
\node[smalllab] at (3.5,4.0) {Dashed links disappear with X};
\node[smalllab] at (3.5,0.0) {The surviving route avoids X};

\end{tikzpicture}
    \end{column}
  \end{columns}
  \assumption{Constructed graph. A targeted route-node failure tests a stress case, not average fleet reliability.}
  \note{A link-disjoint alternative may still share a satellite. Node-disjoint alternatives, apart from endpoints, protect against one internal-node loss only under the graph assumptions. Correlated failures need separate tests.}
\end{frame}

\begin{frame}[fragile]{Compare the Failed Run with Its Matched Baseline}
  \context{Orbital motion continues during a failure, so before-and-after values alone can mix two effects.}
  \begin{columns}[T,onlytextwidth]
    \begin{column}{0.47\textwidth}\raggedright
      \begin{teachingpoints}
        \point{Compare the same timestamps.}{\item At each timestamp, compare both runs under identical motion. Their difference isolates the scheduled failure.}
        \point{Separate outage from extra delay.}{\item Calculate delay differences only where both runs have a route. Report additional unreachable samples separately.}
        \point{Inspect recovery independently.}{\item Restoring a node allows links to return. The model recomputes routes immediately, without protocol recovery delay.}
      \end{teachingpoints}
    \end{column}
    \begin{column}{0.49\textwidth}\centering
      \figtitle{A matched comparison during one event}
{\scriptsize\setlength{\tabcolsep}{4pt}\renewcommand{\arraystretch}{1.35}
\begin{tabular}{@{}llll@{}}\toprule
\textbf{Time} & \textbf{Baseline} & \textbf{Failure} & \textbf{Difference} \\ \midrule
590 s & 12 ms & 12 ms & 0 ms \\
600 s & 13 ms & 18 ms & +5 ms \\
700 s & 14 ms & No path & --- \\
900 s & 15 ms & 15 ms & 0 ms \\
\bottomrule\end{tabular}\par}
\par\vspace{0.35cm}
\begin{tikzpicture}[x=0.9cm,y=0.9cm]
\draw[flow] (0.1,0.55)--(6.7,0.55);
\fill[Mid] (1.5,0.8) rectangle (5.1,1.25);
\node[lab,above] at (3.3,1.35) {Node disabled on $[600,900)$};
\node[smalllab,below] at (1.5,0.5) {600 s};\node[smalllab,below] at (5.1,0.5) {900 s};
\end{tikzpicture}

    \end{column}
  \end{columns}
  \assumption{Constructed comparison values. The companion script supplies actual traces and metrics for the chosen scenario.}
  \note{The interval is half-open. At 600 seconds the satellite is disabled. At 900 it is enabled. A zero difference outside the failure is expected in this deterministic memoryless graph model. It is not a statement about real protocol convergence.}
\end{frame}

\begin{frame}[fragile]{Keep Designs That Offer a Useful Tradeoff}
  \context{No single metric captures satellite count, connectivity, and path quality at the same time.}
  \begin{columns}[T,onlytextwidth]
    \begin{column}{0.47\textwidth}\raggedright
      \begin{teachingpoints}
        \point{Choose a small set of objectives.}{\item Here we prefer fewer satellites, a higher reachable-sample fraction, and lower conditional propagation delay.}
        \point{Reject a dominated design.}{\item B beats C on satellite count and delay while matching reachability. C offers no benefit under these chosen metrics.}
        \point{Keep useful alternatives.}{\item A is cheaper but less connected than B. Choosing between them requires an explicit service goal or resource budget.}
      \end{teachingpoints}
    \end{column}
    \begin{column}{0.49\textwidth}\centering
      \figtitle{An illustrative three-design comparison}
{\scriptsize\setlength{\tabcolsep}{4pt}\renewcommand{\arraystretch}{1.35}
\begin{tabular}{@{}llll@{}}\toprule
\textbf{Design} & \textbf{Satellites} & \textbf{Reachable} & \textbf{Delay} \\ \midrule
A & 48 & $80\%$ & 24 ms \\
B & 72 & $95\%$ & 20 ms \\
C & 96 & $95\%$ & 25 ms \\
\bottomrule\end{tabular}\par}
\par\vspace{0.15cm}
\begin{tikzpicture}
\begin{axis}[width=6.5cm,height=2.75cm,xmin=40,xmax=105,ymin=75,ymax=110,
xtick={48,72,96},ytick={80,95},xlabel={Satellite count},ylabel={Reachable (\%)},
label style={font=\scriptsize},tick label style={font=\scriptsize},grid=major,grid style={Line}]
\addplot[only marks,mark=*,Ink] coordinates {(48,80)(72,95)(96,95)};
\node[font=\scriptsize,anchor=south] at (axis cs:48,81) {A};
\node[font=\scriptsize,anchor=south] at (axis cs:72,96) {B};
\node[font=\scriptsize,anchor=south] at (axis cs:96,96) {C};
\draw[flow] (axis cs:93,92)--(axis cs:75,92);
\end{axis}
\end{tikzpicture}

    \end{column}
  \end{columns}
  \assumption{Constructed results, not simulator output. Satellite count is a resource proxy, not a complete financial cost.}
  \note{Pareto dominance means no worse in every selected objective and strictly better in at least one. The table contains a third objective, delay, that the two-dimensional scatterplot cannot display. Compare delay on a common reachable subset if selection effects matter.}
\end{frame}

\begin{frame}[fragile]{Save the Scenario and the Recorded Outcomes}
  \context{A reproducible result needs more than a screenshot of a favorable instant.}
  \begin{columns}[T,onlytextwidth]
    \begin{column}{0.47\textwidth}\raggedright
      \begin{teachingpoints}
        \point{Preserve inputs and selection rules.}{\item Record the epoch, orbit parameters, topology, endpoints, failure ID, and sampling schedule with the script.}
        \point{Keep the complete recording.}{\item Save states and routes, including no-path samples. A viewer can replay the recording without running Python again.}
        \point{State the model and time window.}{\item Include dependency versions, metrics, and limitations. A short synthetic run does not establish worldwide operational service.}
      \end{teachingpoints}
    \end{column}
    \begin{column}{0.49\textwidth}\centering
      \figtitle{A result another student can inspect}
\begin{tikzpicture}[x=0.90cm,y=0.90cm]
\path[use as bounding box] (-0.25,-0.2) rectangle (7.15,4.5);

\node[box,text width=1.65cm] (a) at (1.0,3.4) {Scenario\\+ script};
\node[box,text width=1.65cm] (b) at (3.4,3.4) {Python\\simulation};
\node[box,text width=1.65cm] (c) at (5.8,3.4) {JSON\\recording};
\draw[flow] (a)--(b);\draw[flow] (b)--(c);
\node[box,text width=2.2cm] (d) at (2.0,1.35) {Metric table\\and assumptions};
\node[box,text width=2.2cm] (e) at (5.35,1.35) {Offline HTML\\viewer};
\draw[flow] (c.south)--(e.north);
\draw[flow] (c.south) to[bend right=20] (d.north);
\node[smalllab,text=Muted] at (3.4,0.15) {The viewer displays records. It does not recompute routes.};

\end{tikzpicture}
    \end{column}
  \end{columns}
  \assumption{The exported trace includes scenario data and dependency versions. The basic model is deterministic.}
  \note{Run.save validates before writing. Run.export\_html packages the existing recording. A recorded seed is not a substitute for specifying a stochastic model, and this basic model draws no random failures.}
\end{frame}

\begin{frame}[fragile]{Build a Defensible Satellite-Network Experiment}
  \context{Combine architecture, coverage, routing, and time-varying behavior in a small final investigation.}
  \begin{columns}[T,onlytextwidth]
    \begin{column}{0.47\textwidth}\raggedright
      \begin{teachingpoints}
        \point{Begin with the paired failure script.}{\item Predict the outcome, run both cases, and locate any extra disconnections or detours in the viewer.}
        \point{Make one justified design revision.}{\item Change the topology policy or constellation layout. Repeat the same failure rule and compare the same metrics and sites.}
        \point{Explain a result and a limitation.}{\item Use a route diagram and a time plot to explain the mechanism. State what the experiment cannot conclude.}
      \end{teachingpoints}
    \end{column}
    \begin{column}{0.49\textwidth}\centering
      \figtitle{A compact final submission}
{\scriptsize\setlength{\tabcolsep}{4pt}\renewcommand{\arraystretch}{1.35}
\begin{tabular}{@{}ll@{}}\toprule
\textbf{Include} & \textbf{Answer} \\ \midrule
Question + prediction & What should change, and why? \\
Baseline + revision & Which design factor changed? \\
Saved JSON + HTML & Can another student inspect it? \\
Metrics + time plot & How often, how long, how far? \\
Conclusion + limits & What does this run support? \\
\bottomrule\end{tabular}\par}
\par\vspace{0.25cm}
\begin{lstlisting}
python teaching/lecture-05/experiment.py
\end{lstlisting}
\par\vspace{0.12cm}{\footnotesize Outputs: paired traces, viewers, and CSV metrics\par}
    \end{column}
  \end{columns}
  \assumption{Use the same failure ID to isolate a design change. Retargeting a new route node asks a different stress-test question.}
  \note{The script saves the chosen failure ID in summary.json. For a revision use that fixed ID when it still identifies a corresponding satellite. If constellation membership changes, define a comparable fault rule explicitly and describe the changed question. Do not claim a throughput improvement from geometric reachability.}
\end{frame}
\end{document}
